\section{Conclusions and Future Work}
\label{sec:discussion}

The experiments of \cref{sec:evaluation} allow the central question of this
thesis to be answered, and the answer is more nuanced than the initial
formulation suggested. This section relates the individual results to that
question, draws the overall conclusions, and outlines where the engine would have
to be extended for the remaining gap to close.

\subsection{Depth and Decision Quality Are Not Proportional}
\label{subsec:discussion_depth_vs_quality}

The clearest finding of this work is that additional search depth does not
translate into proportionally better decisions. Null-move pruning
(\cref{subsec:eval_nmp}) reaches a greater depth at every single one of the nine
time budgets, at one second a full ply, and yet the mean centipawn loss remains
statistically unchanged. Search tree reuse (\cref{subsec:eval_tree_reuse}) adds
another $0.28$ ply at no cost at all. Both refinements do exactly what they are
designed to do, and neither shows up as a measurable gain in move quality on the
test suite.

The move stability measurements (\cref{subsec:eval_move_stability}) explain why.
By depth 8, the depth the engine reaches at a one-second budget, three quarters
of all positions have already settled on their final move, and about a quarter
never change it from the very first iteration. Deeper search can only change the
decision in the remaining fifth of positions, so an across-the-board gain in
depth is diluted accordingly. This is the practical sweet spot the thesis asked
for: it lies where the depth curve is still rising while the quality curve has
already flattened, for this engine and suite at roughly one second per move.

\subsection{Where the Search Is Actually Limited}
\label{subsec:discussion_limits}

The distribution of move quality (\cref{subsec:eval_outliers}) shows where the
remaining errors come from. The engine plays exactly Stockfish's move in
$49\,\%$ of the positions and stays within $25$\,cp in two thirds of them, with a
median loss of $2$\,cp. The mean of $52$\,cp is produced almost entirely by a
small residue: fifteen positions account for nearly half of it.

These positions share a profile. They hold fewer pieces than average, a third of
them are endgames, and three contain a forced mate the engine does not find.
They are precisely the positions in which material and piece-square tables carry
no useful information, since passed pawns, opposition and king activity are not
represented in the evaluation (\cref{subsec:eval_limitations}). In such
positions the limiting factor is not how deep the engine searches but what its
evaluation is able to distinguish, and no attainable depth compensates for that.

Judging stability by evaluation instead of by the identity of the move makes the
consequence explicit (\cref{subsubsec:eval_eps_stability}). While the move
settles in nearly every position, a third of them never come within 20\,cp of the
reference at any depth, ending at an average loss of $120$\,cp. The search does
not remain undecided in these cases, it converges, and it converges on the wrong
move. A stable decision is therefore not by itself a good one, which is worth
keeping in mind wherever move stability is used as a stopping criterion in
time management.

That these positions are limited by evaluation rather than by depth can be shown
directly. They are searched just as deeply as the rest, to the same average of
$10.1$ plies, but they convert that depth into far less: over the course of a
search their loss improves by $5.4$\,cp where the converging positions gain
$36.5$\,cp. Depth is not useless there, it is merely too weak a lever to close a
gap of more than $100$\,cp. This is the clearest argument that further effort
belongs in the evaluation function: a search that optimizes a misleading
objective more thoroughly only becomes more certain of the wrong answer.

This also puts the null result of the parameter sweep
(\cref{subsec:eval_pvs_sweep}) into perspective. Across all 30 configurations the
node count varies by less than two percent, because re-searches are rare to begin
with: even the most aggressive scouting triggers about 104 of them over 1.5
million nodes. The move ordering, built from transposition table moves, MVV-LVA,
killer moves and the history heuristic, is already good enough that the first
move examined is almost always the best one. This is consistent with the
long-standing observation that the effectiveness of Alpha-Beta pruning is
governed by move ordering rather than by the pruning scheme
itself~\cite{knuth1975alphabeta, marsland1986review}, and it means the scouting
parameters have almost no trade-off left to tune.

\subsection{Implications for Time-Managed Search}
\label{subsec:discussion_time}

For iterative deepening under time constraints~\cite{korf1985iterative} these
results are encouraging. The anytime property is not merely a safety net: since
most positions settle early, an interrupted search usually returns the same move
a deeper one would have confirmed. Stopping between iterations therefore costs
far less quality than the depth difference alone would suggest.

The measurements also show that the engine stops being memory-bound well below
the default hash size (\cref{subsec:eval_tt_size}). Only the smallest table
degrades the search measurably, while beyond 16\,MB the curve flattens. Combined
with the tree reuse result, this argues for keeping the table alive across moves
rather than enlarging it: retaining what was already computed helps more than
providing space for more of it, and the benefit grows over the course of a game
as positions transpose more frequently~\cite{warnock1988search}.

Under tournament conditions all of this adds up to a clear result: the engine
wins its match against the reference opponent by roughly $205$ Elo points
(\cref{subsec:engine_eval}), with both colors and a low draw ratio.

\subsection{Limitations}
\label{subsec:discussion_limitations}

Several limitations qualify these conclusions. Move quality is measured against a
strong reference engine rather than against ground truth, so agreement with
Stockfish is treated as a proxy for correctness; positions in which a different
move is equally good are penalised by the agreement metric, which is why
centipawn loss is used as the primary measure.

The suite consists of 309 predefined positions with three repetitions each. It is
large enough for the aggregate effects reported here, but the skewed distribution
of the loss means that individual means are sensitive to a small number of
positions, and differences of a few centipawns should not be over-interpreted.

The opponent in the match was deliberately weakened through Stockfish's built-in
strength limitation. That mechanism provides only a nominal Elo target and works
partly by introducing random inaccuracies, so the reported difference describes
the relation to this specific configuration and not an absolute rating.

Finally, the engine uses a mailbox board representation and a handcrafted
evaluation, both chosen for clarity rather than for strength. Modern engines rely
on bitboards~\cite{cpw_bitboards} and learned evaluation
functions~\cite{silver2018alphazero}, and the conclusions about absolute playing
strength do not transfer to those designs.

\subsection{Future Work}
\label{subsec:discussion_future}

The results point to one clear priority. Since the residual errors are
concentrated in positions the evaluation cannot describe, further work should be
directed at the evaluation function rather than at the search. Pawn structure,
passed pawns, king activity and dedicated endgame knowledge would address exactly
the positions that dominate the mean loss today, and only then would the depth
already gained by null-move pruning and tree reuse have anything left to
contribute.

A second direction concerns the measurement itself. Reporting the median
alongside the mean, or trimming the extreme cases, would make the effect of
search refinements visible at all, since the mean is dominated by a residue that
no search improvement can fix.

The stability analysis could be extended in the same spirit. Judging convergence
by evaluation rather than by the identity of the move already proved more
informative here (\cref{subsubsec:eval_eps_stability}); applying the same
criterion to the time-to-quality measurements, and reporting from which
\emph{budget} rather than from which depth the loss stays within a margin, would
describe the practical behavior of the engine under a clock more directly than
the depth-based view can.

Finally, the search refinements evaluated here are only the beginning of what a
classical engine typically employs. Late move reductions, aspiration windows and
static exchange evaluation all target the same goal of spending effort where it
changes the decision, and the move ordering measured here is strong enough that
they would likely pay off.

\subsection{Summary}

This thesis implemented a complete chess engine and used it to examine how
invested search time translates into decision quality. Both refinements studied
work as intended: null-move pruning gains a full ply at one second, and a
persistent transposition table adds a further $0.28$ ply for free, growing to
$0.39$ ply in the endgame. Neither improves the average move quality measurably,
because three quarters of all positions have settled by the depth the engine
already reaches, and the remaining errors are concentrated in endgames that the
evaluation cannot assess regardless of depth.

The sweet spot between search time and decision quality is therefore reached
earlier than expected, at around one second per move for this engine. Beyond that
point, additional time keeps buying depth but no longer buys decisions. What
limits the engine at that stage is no longer how far it looks ahead, but what it
is able to recognise when it gets there.
